\subsection{Coverage and independent boundary controls}

On 84 source inputs with the existing two-million-guard allowance, the
four-tree policy now has 24 native positives, 59 completed restricted
misses and one cap, compared with 23 positives before annulus recognition.
Requesting 24 tree trials now gives 27 positives, 56 misses and one cap,
compared with 24 positives before. The capped source remains Gordian.
The optional four-face-root policy retains its 24 positives; its newly
found disc is now preempted by the earlier optimized annulus.

Three earlier positives, \code{random-29}, \code{random-79} and
\code{random-99}, now stop at their raw annuli. Their shared guard work
falls from 11,227, 10,192 and 7,785 to 6,561, 6,480 and 4,984 respectively.
The seven-crossing \code{random-142} stops on its original 122-piece
optimized annulus at 93,012 guards. Without annulus recognition and without
face search it previously completed a miss at 93,011 guards: the extra
unit is the final success guard, not another discovery pass. With four
face roots enabled, its old disc proof required 117,848 guards.

With the optional 24-tree allowance, \code{random-81} now succeeds at
trial 5 rather than trial 22, reducing guards from 72,527 to 52,641.
Two additional four-crossing unknots, \code{random-1} and
\code{random-127}, succeed at tree trials 10 and 7 respectively. Their
annuli have 86 and 111 pieces, and their guard totals are 44,244 and 48,849.
These are additional positives for that explicit extended allowance;
they remain misses under the default four tree trials.

Seven distinct annulus certificates are retained. Every prior positive
certificate still passes the public verifier. Under each of the three
compared policies, twenty positive records retain their exact certificate
hash; the intentional changes use the new annulus tag and may stop earlier.
Cases that do not return an annulus preserve their complete recorded native
result, including their work count. On nineteen ordinary corpus inputs,
both settings return the expected verdicts without reaching the native
stage, so no ordinary-portfolio gain follows from this experiment.

The final audit independently checks all 787 original or root-extremal
minimum-span surfaces from the preceding 74-source small-input experiment.
Sixteen have Euler characteristic zero, and precisely those pass the new
annulus wrapper. Regina agrees on connectedness, orientability, Euler
characteristic and boundary count. The seven actual produced annulus
witnesses receive separate Regina comparisons, including raw and later-tree
surfaces outside that minimum-span experiment. Two trefoil adversaries
bring the total to 796 Regina comparisons.

For the sixteen accepted minimum-span annuli and the seven produced
witnesses, a separate small-fixture control expands boundary arc pairings,
finds each boundary circle and computes its mod-two intersection vector
against an independently validated torus homology basis. All 23 comparisons
give two circles, exactly one with zero vector. At most 86 boundary points
are expanded. This is audit code only; neither recognition nor certificate
verification calls it.

The first trefoil adversary is a genuine connected orientable annulus with
20 normal pieces. Set the local height at a corner to one exactly when its
canonical global vertex is 3 or 7, and to zero otherwise. This is a pure
vertex coboundary. Its two boundary circles both have zero mod-two homology,
and the primitive-class checker rejects it. The second adversary changes
the primitive trefoil seed by subtracting one at global vertex 3. Its total
Euler characteristic is zero, but Regina and the native component census
find two components and no compressing disc. The zero-tree connectedness
check rejects this forged annulus witness.

\subsection{Tests and complete-call measurements}

All 34 focused tests pass in 4.809 seconds, and all 1,204 maintained tests
pass in 226.376 seconds. Tests cover producer-free annulus replay, raw and
optimized early termination, exact budget boundaries, cancellation, strict
switch types, schema swaps, malformed fields, wrong sources, the trefoil
adversaries, zero and doubled classes, and 20,000-bit gauge transport through
hexadecimal JSON. Historical tree and face tests explicitly disable the new
criterion when testing their original discovery paths.

The final audit takes 159.502 seconds and pins 488 source files. Its
elapsed time is a correctness-run total, not a paired performance result.
A preliminary successful audit is retained as a log; the final run follows
an update that makes the historical tree, blocking-flow and face research
drivers explicitly retain their original disc-only comparison criterion.
The cocycle-surface inspector and span verifier remain byte-identical to
baseline \code{b2d13b5f4}.

Measurements run serially without overlapping tests, audits, builds or
commits. They start from the input diagram and include the complete
recognizer, exact fallback when needed, independent replay of any returned
native proof and certificate hashing. Earlier shortcuts are disabled to
reach the native stage. Both arms have no local seed cap and retain the
twelve-second global deadline. The on/off comparison changes only the
annulus criterion. Tree trials are four except for the three explicitly
extended-tree cases; face roots are zero except for the labeled prior
face-search positive, which requests four.

The initial five-round run has visible A/A variation, including a 1.089
raw-disc on-arm ratio and a 0.876 ratio for the three-crossing raw annulus.
It is retained, and a nine-round repeat supplies the final table. Each run
has one warm-up round per case, and each round randomizes off, off A/A,
on and on A/A endpoints. All 616 measured calls and 88 warm-ups across the
two runs finish. Their per-implementation verdicts, witness hashes and
native work counts are stable across both runs. The final run has 396
measured calls and 44 warm-ups and takes 37.153 seconds.

\begin{center}
\small\begin{tabular}{lrrrrr}
\toprule
Input & Off (ms) & On (ms) & Off/on & Off A/A & On A/A \\
\midrule
Raw disc control & 64.707 & 65.940 & 0.992 & 0.958 & 0.990 \\
Optimized disc control & 91.996 & 90.158 & 1.012 & 1.008 & 0.995 \\
Raw annulus, 4 crossings A & 50.989 & 31.007 & 1.637 & 1.001 & 0.995 \\
Raw annulus, 4 crossings B & 49.219 & 31.143 & 1.580 & 0.998 & 0.999 \\
Raw annulus, 3 crossings & 36.943 & 23.334 & 1.571 & 0.997 & 1.008 \\
New optimized annulus & 118.295 & 137.411 & 0.865 & 1.005 & 0.991 \\
Prior face-search positive & 206.528 & 137.443 & 1.483 & 0.998 & 0.988 \\
Earlier late-tree positive & 136.703 & 100.128 & 1.349 & 1.031 & 0.970 \\
New tree annulus A & 89.212 & 85.510 & 1.033 & 1.021 & 1.004 \\
New tree annulus B & 92.456 & 85.370 & 1.075 & 0.993 & 0.998 \\
Trefoil miss & 39.753 & 40.171 & 0.992 & 0.998 & 1.001 \\
\bottomrule
\end{tabular}

\end{center}

The three early-annulus cases improve $1.57$--$1.64\times$ by paired
ratios. The previous face-search positive improves $1.48\times$, and the
late-tree positive improves $1.35\times$. The two new extended-tree
positives have ratios 1.033 and 1.075; the first is modest relative to its
A/A variation, so the main conclusion there is added native coverage.
The unchanged raw-disc, optimized-disc and trefoil controls are near parity.
Table times are primary-arm medians and ratios are medians of paired
ratios; these need not be quotients of the displayed median times.

The newly recognized seven-crossing source without face search is an
adverse complete-call control. Its old exact fallback finishes in a median
118 milliseconds; returning and independently replaying the new native
proof takes 137 milliseconds, about 16 percent slower by the paired ratio.
There is no claim of a uniform recognition speedup. The criterion adds no
surface discovery loop and shortens several existing native searches, but
native proof production and replay can cost more than the old fallback on
a small diagram.

The reproduction driver is \path{../fast/normal_orbit_research/annulus.py}.
Use its \code{audit} or \code{recognize --rounds 9} subcommand;
the audit uses optional Regina.
The incoming archives and checked \code{openai/math} head
\code{fd4aeeb2e} are unchanged. This is a stronger polynomial sufficient
certificate test within an incomplete candidate search, not a general
subexponential or quasipolynomial recognition algorithm.
